\documentclass[10pt,a4paper]{amsart}
\usepackage{amsmath,amssymb,amsthm,graphicx}
% Retained preamble for the Stacks Project chapter morphisms.tex (Morphisms of Schemes).
% The chapter file pinned-original-source.tex carries only \input{preamble} on its first line, so no
% preamble text lives inside the chapter itself.  The Stacks Project uses one single shared preamble for
% every chapter; the snapshot stages it verbatim in sources/stacks-chow/preamble.tex.  The block below is
% copied mechanically, line for line, from that file's lines 14-18 (the Xy-pic setup) and lines 154-217
% (the native theorem declarations together with the mathematical macros).  Nothing was added, removed or
% redefined.  The wrapper class and the amsmath/amssymb/amsthm/graphicx base packages are supplied by the
% assembler, so the class/\externaldocument part of the shared preamble is not repeated here.
% For commutative diagrams we use Xy-pic
\usepackage[all]{xy}

% We use 2cell for 2-commutative diagrams.
\xyoption{2cell}
% Theorem environments.
%
\theoremstyle{plain}
\newtheorem{theorem}[subsection]{Theorem}
\newtheorem{proposition}[subsection]{Proposition}
\newtheorem{lemma}[subsection]{Lemma}

\theoremstyle{definition}
\newtheorem{definition}[subsection]{Definition}
\newtheorem{example}[subsection]{Example}
\newtheorem{exercise}[subsection]{Exercise}
\newtheorem{situation}[subsection]{Situation}

\theoremstyle{remark}
\newtheorem{remark}[subsection]{Remark}
\newtheorem{remarks}[subsection]{Remarks}

\numberwithin{equation}{subsection}

% Macros
%
\def\lim{\mathop{\mathrm{lim}}\nolimits}
\def\colim{\mathop{\mathrm{colim}}\nolimits}
\def\Spec{\mathop{\mathrm{Spec}}}
\def\Hom{\mathop{\mathrm{Hom}}\nolimits}
\def\Ext{\mathop{\mathrm{Ext}}\nolimits}
\def\SheafHom{\mathop{\mathcal{H}\!\mathit{om}}\nolimits}
\def\SheafExt{\mathop{\mathcal{E}\!\mathit{xt}}\nolimits}
\def\Sch{\mathit{Sch}}
\def\Mor{\mathop{\mathrm{Mor}}\nolimits}
\def\Ob{\mathop{\mathrm{Ob}}\nolimits}
\def\Sh{\mathop{\mathit{Sh}}\nolimits}
\def\NL{\mathop{N\!L}\nolimits}
\def\CH{\mathop{\mathrm{CH}}\nolimits}
\def\proetale{{pro\text{-}\acute{e}tale}}
\def\etale{{\acute{e}tale}}
\def\QCoh{\mathit{QCoh}}
\def\Ker{\mathop{\mathrm{Ker}}}
\def\Im{\mathop{\mathrm{Im}}}
\def\Coker{\mathop{\mathrm{Coker}}}
\def\Coim{\mathop{\mathrm{Coim}}}

% Boxtimes
%
\DeclareMathSymbol{\boxtimes}{\mathbin}{AMSa}{"02}

%
% Macros for moduli stacks/spaces
%
\def\QCohstack{\mathcal{QC}\!\mathit{oh}}
\def\Cohstack{\mathcal{C}\!\mathit{oh}}
\def\Spacesstack{\mathcal{S}\!\mathit{paces}}
\def\Quotfunctor{\mathrm{Quot}}
\def\Hilbfunctor{\mathrm{Hilb}}
\def\Curvesstack{\mathcal{C}\!\mathit{urves}}
\def\Polarizedstack{\mathcal{P}\!\mathit{olarized}}
\def\Complexesstack{\mathcal{C}\!\mathit{omplexes}}
% \Pic is the operator that assigns to X its picard group, usage \Pic(X)
% \Picardstack_{X/B} denotes the Picard stack of X over B
% \Picardfunctor_{X/B} denotes the Picard functor of X over B
\def\Pic{\mathop{\mathrm{Pic}}\nolimits}
\def\Picardstack{\mathcal{P}\!\mathit{ic}}
\def\Picardfunctor{\mathrm{Pic}}
\def\Deformationcategory{\mathcal{D}\!\mathit{ef}}

% --- appended 2026-10-12 (Chief): the project own definitions of the reference/slogan/history
% display environments, copied mechanically from the shared Stacks preamble lines 8-12
% (\usepackage{verbatim} plus the three \newenvironment lines, which the chapter preamble keeps
% verbatim from the project preamble). Nothing added, removed or redefined: these environments
% carry no mathematical content and the project itself typesets them as comments. Without these
% three lines a counted statement whose own text uses them fails with Environment slogan undefined.
\usepackage{verbatim}
\newenvironment{reference}{\comment}{\endcomment}
\newenvironment{slogan}{\comment}{\endcomment}
\newenvironment{history}{\comment}{\endcomment}
\usepackage{cleveref}
\crefname{originallemma}{original lemma}{original lemmas}
\Crefname{originallemma}{Original lemma}{Original lemmas}
\makeatletter
\newlabel{morphisms-lemma-integral-closure}{{035F}{source}{}{sourcepointer-1}{}}
\newlabel{morphisms-lemma-integral-closure@cref}{{[originallemma][0][]035F}{[1][0][]source}{}{}{}}
\makeatother
\crefname{originallemma}{original lemma}{original lemmas}
\Crefname{originallemma}{Original lemma}{Original lemmas}
\makeatletter
\newlabel{algebra-lemma-integral-closure-is-ring}{{00GO}{source}{}{sourcepointer-2}{}}
\newlabel{algebra-lemma-integral-closure-is-ring@cref}{{[originallemma][0][]00GO}{[1][0][]source}{}{}{}}
\makeatother
\crefname{originallemma}{original lemma}{original lemmas}
\Crefname{originallemma}{Original lemma}{Original lemmas}
\makeatletter
\newlabel{algebra-lemma-integral-closure-stalks}{{034K}{source}{}{sourcepointer-3}{}}
\newlabel{algebra-lemma-integral-closure-stalks@cref}{{[originallemma][0][]034K}{[1][0][]source}{}{}{}}
\makeatother
\crefname{originallemma}{original lemma}{original lemmas}
\Crefname{originallemma}{Original lemma}{Original lemmas}
\makeatletter
\newlabel{algebra-lemma-integral-closure-localize}{{0307}{source}{}{sourcepointer-4}{}}
\newlabel{algebra-lemma-integral-closure-localize@cref}{{[originallemma][0][]0307}{[1][0][]source}{}{}{}}
\makeatother
\title{The sheaf of elements integral over the structure sheaf}
\author{The Stacks Project Authors}
\date{stacks-project:morphisms.tex:local-snapshot-20261012; GNU FDL 1.2 or later}
\begin{document}
\maketitle
\setcounter{subsection}{1}
\noindent\textit{Editorial note.} All delivered mathematical text is the author's own, copied from the pinned Stacks Project chapter file morphisms.tex; no statement or proof step was rewritten, reordered or removed.\par\medskip
\noindent\textit{Editorial note.} The chapter file pinned-original-source.tex carries only \textbackslash{}input\{preamble\} on its first line, so no preamble text lives inside the chapter. The retained file sources/stacks-morphisms/retained-preamble.tex is a mechanical, line-for-line copy of the project's single shared preamble as staged in the snapshot. Nothing was added, removed or redefined, and the excerpt serves only as preamble, never as a span coordinate.\par\medskip
\noindent\textit{Editorial note.} Both counted spans start at the \textbackslash{}begin line of their environment, so the statement and the author's proof are delivered whole; surrounding prose of the section is deliberately outside the counted spans and is carried as standard prerequisites instead.\par\medskip
\noindent\textit{Editorial note.} The citation fields (source\_id, source\_version, authors, source\_url, license.name, source\_title) are identical to those of every other card this lane has produced from this chapter, so the batch selector treats them as one source.\par\medskip
\noindent\textit{Allowed standard prerequisites (editorial).}\begin{itemize}
\item Integral elements of a ring extension form a subring. Algebra, Lemma \textbackslash{}ref\{algebra-lemma-integral-closure-is-ring\}, cited at proof lines 13862-13864. Makes the collection a sheaf of algebras.
\item Description of the integral closure through stalks in the affine case. Algebra, Lemma \textbackslash{}ref\{algebra-lemma-integral-closure-stalks\}, cited at proof lines 13865-13870. Gives the stalk and affine-open identifications.
\end{itemize}\medskip
\ifdefined\hypertarget\hypertarget{sourcepointer-1}{}\fi
\noindent\textit{Source reference.} Original lemma 035F: focused-seeds/AGP/sources/stacks-snapshot-20261011/tags.txt, source lines 4126--4126. The lemma's own canonical tag, quoted so that readers can resolve the label inside the delivered text; it points at the counted statement itself and carries no proof content.\par\medskip
\ifdefined\hypertarget\hypertarget{sourcepointer-2}{}\fi
\noindent\textit{Source reference.} Original lemma 00GO: focused-seeds/AGP/sources/stacks-snapshot-20261011/tags.txt, source lines 614--614. Navigation pointer to a statement the counted proof refers to but does not contain; nothing load-bearing is delegated to it.\par\medskip
\ifdefined\hypertarget\hypertarget{sourcepointer-3}{}\fi
\noindent\textit{Source reference.} Original lemma 034K: focused-seeds/AGP/sources/stacks-snapshot-20261011/tags.txt, source lines 4095--4095. Navigation pointer to a statement the counted proof refers to but does not contain; nothing load-bearing is delegated to it.\par\medskip
\ifdefined\hypertarget\hypertarget{sourcepointer-4}{}\fi
\noindent\textit{Source reference.} Original lemma 0307: focused-seeds/AGP/sources/stacks-snapshot-20261011/tags.txt, source lines 3938--3938. Navigation pointer to a statement the counted proof refers to but does not contain; nothing load-bearing is delegated to it.\par\medskip
\begin{lemma}
\label{lemma-integral-closure}
Let $X$ be a scheme. Let $\mathcal{A}$ be a quasi-coherent sheaf
of $\mathcal{O}_X$-algebras. The subsheaf $\mathcal{A}' \subset \mathcal{A}$
defined by the rule
$$
U \longmapsto \{f \in \mathcal{A}(U) \mid
f_x \in \mathcal{A}_x \text{ integral over } \mathcal{O}_{X, x}
\text{ for all }x \in U\}
$$
is a quasi-coherent $\mathcal{O}_X$-algebra, the stalk $\mathcal{A}'_x$
is the integral closure of $\mathcal{O}_{X, x}$ in $\mathcal{A}_x$, and
for any affine open $U \subset X$ the ring
$\mathcal{A}'(U) \subset \mathcal{A}(U)$ is
the integral closure of $\mathcal{O}_X(U)$ in $\mathcal{A}(U)$.
\end{lemma}

\begin{proof}
This is a subsheaf by the local nature of the conditions.
It is an $\mathcal{O}_X$-algebra by
Algebra, Lemma \ref{algebra-lemma-integral-closure-is-ring}.
Let $U \subset X$ be an affine open. Say $U = \Spec(R)$
and say $\mathcal{A}$ is the quasi-coherent sheaf associated to
the $R$-algebra $A$. Then according to
Algebra, Lemma \ref{algebra-lemma-integral-closure-stalks}
the value of $\mathcal{A}'$ over $U$ is given by the integral
closure $A'$ of $R$ in $A$. This proves the last assertion of
the lemma. To prove that $\mathcal{A}'$ is quasi-coherent, it
suffices to show that $\mathcal{A}'(D(f)) = A'_f$. This follows
from the fact that integral closure and localization commute, see
Algebra, Lemma \ref{algebra-lemma-integral-closure-localize}.
The same fact shows that the stalks are as advertised.
\end{proof}

\end{document}
